% Compile with pdflatex ExtraPractice.tex from this directory.
% To use the cached Windows compiler, run these lines from the repository root:
% $practiceRoot = (Resolve-Path '.\midterms\midterm-2-practice').Path
% $day4 = Join-Path $practiceRoot 'practice/day-04-dafny'
% $env:TECTONIC_CACHE_DIR = Join-Path $practiceRoot '.tools/cache'
% $env:FONTCONFIG_FILE = Join-Path $practiceRoot '.tools/fonts.conf'
% & "$practiceRoot\.tools\tectonic.exe" --only-cached --untrusted --outdir $day4 "$day4\ExtraPractice.tex"
\documentclass[11pt,letterpaper]{article}
\usepackage[margin=0.8in]{geometry}
\usepackage[T1]{fontenc}
\usepackage{lmodern}
\usepackage{amsmath,xcolor,listings,hyperref}

\hypersetup{pdftitle={CMSC433 Day 4 - Extra Dafny practice},
  pdfsubject={Additional practice in weakest preconditions and loop obligations}}
\definecolor{ink}{HTML}{172E3B}
\definecolor{accent}{HTML}{1E586B}
\definecolor{paper}{HTML}{F5F8FA}
\definecolor{line}{HTML}{718693}
\color{ink}
\setlength{\parindent}{0pt}
\setlength{\parskip}{0.38em}
\raggedbottom

\lstset{basicstyle=\small\ttfamily,columns=fullflexible,keepspaces=true,
  showstringspaces=false,backgroundcolor=\color{paper},frame=single,
  rulecolor=\color{line},framesep=6pt,xleftmargin=7pt,xrightmargin=7pt,
  aboveskip=0.65em,belowskip=0.65em}

\newcommand{\question}[2]{%
  \par{\Large\bfseries\color{accent}#1\enspace #2}\par\smallskip}

% Each \answer takes a field ID, a label, a PDF field height, and your answer.
% To edit the LaTeX, type in the final braces. An empty answer becomes a
% fillable field. A written answer is typeset as part of the page on rebuild.
% Use \(...\) for math and change the height if you want more PDF writing space.
\renewcommand*{\LayoutTextField}[2]{#1#2}
\newcommand{\answer}[4]{%
  \par\smallskip\noindent\textbf{#2}\par\nobreak\noindent
  \if\relax\detokenize{#4}\relax
    % Draw the field's outline on the page for print and PDF previews.
    \rlap{\setlength{\fboxsep}{0pt}%
      \raisebox{\fboxrule}[0pt][0pt]{\fcolorbox{line}{paper}{%
        \rule{0pt}{\dimexpr#3-2\fboxrule\relax}%
        \rule{\dimexpr\linewidth-2\fboxrule\relax}{0pt}}}}%
    \TextField[name={#1},multiline=true,width=\linewidth,height=#3,
      charsize=11pt,borderwidth=1,bordercolor={0.44 0.53 0.58},
      backgroundcolor={0.96 0.97 0.98}]{}%
  \else
    \fcolorbox{line}{paper}{%
      \parbox{\dimexpr\linewidth-2\fboxsep-2\fboxrule\relax}{#4}}%
  \fi
  \par}

\begin{document}
\begin{Form}

{\large\bfseries\color{accent} CMSC433 / Day 4 / Extra practice}\par
{\LARGE\bfseries Dafny verification conditions}\par\smallskip
Work one step at a time. Type in the PDF fields or put your answers in the
final braces of each \texttt{\string\answer} in the LaTeX source.
Use partial correctness throughout.

\subsection*{Worked example}
To get \(x\geq 7\) after \texttt{x := x + 3}, you need \(x+3\geq 7\)
before the assignment. Replace the assigned variable in the target with the
right-hand side of the assignment. For a sequence, start with the last
statement and work backward.

\smallskip
{\small
\renewcommand{\arraystretch}{1.23}
\begin{tabular}{@{}p{.20\linewidth}p{.34\linewidth}p{.39\linewidth}@{}}
\textbf{Loop obligation} & \textbf{Assumptions} & \textbf{Condition to establish}\\
\hline
Initialization & Method precondition & Invariant after the prefix\\
Preservation & Invariant and true guard & Invariant after one body execution\\
Exit & Invariant and false guard & Method postcondition after the suffix\\
\end{tabular}\par}
\medskip

\question{4.1}{Weakest precondition}
Work backward through these assignments. Keep each substitution visible.

\begin{lstlisting}
x := y - 3;
y := 2*x;
// Target: x < y && y >= 4
\end{lstlisting}

\answer{q41-before-second}{Before the second assignment}{0.9in}{}
\answer{q41-before-first}{Before the first assignment}{0.9in}{}

\medskip
{\small Course review 27:47--38:06, Dafny verification lecture, and
Spring 2025 sample exam Question 3.}\par

\clearpage
\question{4.2}{Initialization, preservation, and exit}
Use the invariant \(I\equiv i\leq n\land t=4i\).
Derive one implication at a time. Use the method precondition in the
initialization implication and account for the suffix in the exit implication.

\begin{lstlisting}
method CountUpByFour(n: int) returns (t: int)
  requires n >= 1
  ensures t == 4*n - 2
{
  t := 4;
  var i := 1;
  while i < n {
    t := t + 4;
    i := i + 1;
  }
  t := t - 2;
}
\end{lstlisting}

\answer{q42-initialization}{Initialization implication}{0.95in}{}
\answer{q42-preservation}{Preservation implication}{0.95in}{}
\answer{q42-exit}{Exit implication}{0.95in}{}

\clearpage
\question{4.3}{Compare proposed invariants}
Use the \texttt{CountUpByFour} program from 4.2. For each candidate, derive all
three implications. Identify each failed implication and give concrete values
that make its left-hand side true and its right-hand side false.

\medskip
{\bfseries\color{accent} Candidate A}\quad \(i\leq n\land t=4n\)\par
\answer{q43-a-initialization}{Initialization implication}{0.52in}{}
\answer{q43-a-preservation}{Preservation implication}{0.52in}{}
\answer{q43-a-exit}{Exit implication}{0.52in}{}
\answer{q43-a-counterexamples}{Failed implications and counterexamples}{0.65in}{}

\medskip
{\bfseries\color{accent} Candidate B}\quad \(t=4i\)\par
\answer{q43-b-initialization}{Initialization implication}{0.52in}{}
\answer{q43-b-preservation}{Preservation implication}{0.52in}{}
\answer{q43-b-exit}{Exit implication}{0.52in}{}
\answer{q43-b-counterexamples}{Failed implications and counterexamples}{0.65in}{}

\clearpage
\question{4.4}{Transfer}
Use the original invariant \(I\equiv i\leq n\land t=4i\) from 4.2.
Replace the suffix and postcondition with the following.

\begin{lstlisting}
t := t + 3;
t := 2*t;
// Target: t == 8*n + 6
\end{lstlisting}

Derive the new exit implication and show both substitutions. Explain why
initialization and preservation stay the same.

\answer{q44-before-last}{Before the last assignment in the suffix}{0.8in}{}
\answer{q44-before-first}{Before the first assignment in the suffix}{0.8in}{}
\answer{q44-exit}{New exit implication}{0.9in}{}
\answer{q44-explanation}{Why initialization and preservation stay the same}{1.15in}{}

\end{Form}
\end{document}
